        %auto-ignore
\section{Related Works}
\subsection{Width Scaling and $\mu$P} The infinite-width limit of neural networks has been a topic of extensive research in the literature. Numerous studies have predominantly focused on examining the behavior of various statistical quantities at initialization. Some works have gone beyond the initialization stage to explore the dynamics of feature learning in neural networks.
\paragraph{Lazy training.} With standard parametrization, a learning rate of order $\mathcal{O}(n^{-1})$,\footnote{We also obtain the lazy infinite-width limit with the NTK parametrization and a $\mathcal{O}(n^{-1/2})$ learning rate.} $n$ being the width, yields the so-called lazy training regime in the infinite-width limit, where the features remain roughly constant throughout training \citep{chizat2020lazy, yang2022tensor}. This regime is also known as the Neural Tangent Kernel (NTK) regime and its convergence properties have been extensively studied in the literature \citep{jacot2020neural, allenzhu2019convergence, chizat2018global, zou2018stochastic}. 

\paragraph{Feature learning and $\mu$P.} Recent empirical studies (e.g. \cite{yang2022tensor}) have provided compelling evidence that feature learning plays a crucial role in the success of deep learning. It is widely acknowledged that the remarkable performance achieved by deep neural networks can be attributed to their ability to acquire meaningful representations through the process of training. Consequently, scaling the network architecture emerges as a natural choice to enhance the performance of such models.

In this context, $\mu$P (Maximal Update Parameterization), introduced in \cite{yang2022tensor}, has emerged as a promising approach for maximizing feature learning while simultaneously preventing feature explosion as the network width increases, given a fixed depth. Notably, $\mu$P facilitates hyperparameter transfer across varying network widths. This means that instead of tuning hyperparameters directly on large models, one can optimize them on smaller models and utilize the same set of hyperparameters for larger models.

The derivation of $\mu$P leverages the Tensor Programs framework \citep{yang2021tensor_i, yang2020scaling, yang2020tensor, yang2021tensor, yang2022tensor}, which provides valuable tools for capturing the behavior of neural networks in the infinite-width regime during the training process. 

\subsection{Depth Scaling} 
While increasing the width of neural networks can lead to improved performance, increasing the depth of the network also yields significant performance gains, and most state-of-the-art models use deep architectures. The introduction of skip connections \citep{he2016deep, he2016identity} played a pivotal role in enabling the training of deep networks. However, it became apparent that even with skip connections and normalization layers, training deep networks remains a challenging task \citep{liu2020understanding}. Moreover, tuning hyperparameters for large depth networks is a time-and-resource-consuming task.

To address the challenges associated with training deep networks, several studies have proposed scaling the network blocks using a depth-dependent scaler to ensure stability of features and gradients at initialization or in the kernel regime \citep{hayou21stable, hanin2018start, zhang2019fixup, noci2022signal, hayou2023infinitedepth, hayou2023width, noci2023shaped,zhang2023stabilize}. However, these works lack insights into the dynamics with feature learning. For instance, one might argue that features can still experience explosive growth if the learning rate is not properly chosen. Therefore, an effective depth scaling approach should not only ensure stability at initialization but also provide guidelines for scaling the learning rate.

This motivation underlies the development of Depth-$\mu$P, which offers a comprehensive framework for depth scaling. Depth-$\mu$P encompasses block multipliers and learning rate scaling, providing a complete recipe for training deep networks. In the case of Multi-Layer Perceptrons (MLPs) (no skip connections), \citet{jelassi2023depth} showed that a learning rate scaling of $depth^{-3/2}$ guarantees stability after the initial gradient step. However, it remains unclear how the learning rate should be adjusted beyond the first step, and this scaling is not suitable for architectures with residual connections.

